%==============================================================================
% Chapitre 37 - La formule de Miller
%==============================================================================

\section{Introduction}

La formule de Greenhill constitue une méthode simple permettant d'estimer le
pas de rayure nécessaire à la stabilisation d'un projectile.

Cependant, elle présente plusieurs limitations et ne fournit pas directement
une estimation du facteur de stabilité gyroscopique.

Afin d'améliorer cette approche, Don Miller a proposé au début du XXI\ieme\
siècle une nouvelle méthode mieux adaptée aux projectiles modernes.

\begin{aretenir}

La formule de Miller permet d'estimer directement le facteur de stabilité
gyroscopique d'un projectile.

\end{aretenir}

\section{Objectif de la formule}

La formule de Miller vise à répondre à la question suivante :

\begin{quote}
Le projectile est-il suffisamment stabilisé par le canon utilisé ?
\end{quote}

La réponse est obtenue sous la forme d'un facteur de stabilité :

\begin{equation}
S_g
\end{equation}

dont l'interprétation a été présentée au chapitre précédent.

\section{Principe général}

La formule prend en compte plusieurs paramètres :

\begin{itemize}
    \item la masse du projectile ;
    \item son diamètre ;
    \item sa longueur ;
    \item le pas de rayure ;
    \item les conditions atmosphériques.
\end{itemize}

Elle est donc plus représentative que la formule de Greenhill.

\section{Version simplifiée}

Une forme simplifiée de la formule peut s'écrire :

\begin{equation}
S_g=
\frac{30m}
     {d^3l(1+l^2)}
\left(\frac{t}{10}\right)^2
\left(\frac{V}{2800}\right)^{1/3}
\end{equation}

où :

\begin{itemize}
    \item $m$ représente la masse du projectile ;
    \item $d$ représente le diamètre ;
    \item $l$ représente la longueur exprimée en calibres ;
    \item $t$ représente le pas de rayure ;
    \item $V$ représente la vitesse initiale.
\end{itemize}

\section{Remarque importante}

La formule existe sous plusieurs variantes.

Les logiciels balistiques modernes utilisent souvent des versions modifiées
tenant compte :

\begin{itemize}
    \item de la densité de l'air ;
    \item de la température ;
    \item de la pression atmosphérique ;
    \item de la forme du projectile.
\end{itemize}

Les équations exactes peuvent donc varier légèrement selon les sources.

\section{Influence de la masse}

La stabilité augmente généralement avec la masse du projectile.

Toutefois, la masse n'est pas le paramètre dominant.

Deux projectiles de même masse peuvent présenter des stabilités très
différentes.

\section{Influence de la longueur}

La longueur joue un rôle majeur.

Plus un projectile est long :

\begin{itemize}
    \item plus il est difficile à stabiliser ;
    \item plus il nécessite une rotation importante.
\end{itemize}

Cette observation rejoint les conclusions obtenues avec Greenhill.

\begin{aretenir}

La longueur du projectile constitue l'un des paramètres les plus influents
pour la stabilité.

\end{aretenir}

\section{Influence du diamètre}

Le diamètre intervient également dans le calcul.

À longueur égale :

\begin{itemize}
    \item un projectile de plus gros diamètre est généralement plus facile à
          stabiliser ;
    \item un projectile de faible diamètre nécessite davantage de rotation.
\end{itemize}

\section{Influence du pas de rayure}

Le pas de rayure agit directement sur la vitesse de rotation.

Un pas plus court :

\begin{itemize}
    \item augmente la vitesse de rotation ;
    \item augmente généralement la stabilité gyroscopique.
\end{itemize}

Cette relation explique l'intérêt de choisir un pas adapté au projectile.

\section{Influence de la vitesse}

La vitesse initiale intervient également.

Lorsque la vitesse augmente :

\begin{itemize}
    \item la rotation augmente ;
    \item certains effets aérodynamiques évoluent ;
    \item la stabilité peut être modifiée.
\end{itemize}

\section{Influence de l'atmosphère}

L'atmosphère joue un rôle important.

La stabilité dépend notamment :

\begin{itemize}
    \item de la température ;
    \item de la pression ;
    \item de l'altitude ;
    \item de la densité de l'air.
\end{itemize}

C'est pourquoi un projectile peut se comporter différemment selon le lieu de
tir.

\section{Exemple qualitatif}

Considérons deux projectiles de même calibre :

\begin{itemize}
    \item projectile A : court ;
    \item projectile B : long.
\end{itemize}

À vitesse et pas de rayure identiques :

\begin{itemize}
    \item le projectile A présentera généralement un facteur de stabilité
          plus élevé ;
    \item le projectile B sera plus difficile à stabiliser.
\end{itemize}

\section{Interprétation du facteur de stabilité}

Les valeurs suivantes sont souvent utilisées comme repères :

\begin{itemize}
    \item $S_g < 1$ : instable ;
    \item $S_g \approx 1$ : stabilité marginale ;
    \item $1,3 < S_g < 2$ : zone généralement recherchée ;
    \item $S_g > 2$ : stabilité élevée.
\end{itemize}

Ces limites doivent être considérées comme des ordres de grandeur.

\section{Avantages de la formule}

Par rapport à Greenhill, la formule de Miller :

\begin{itemize}
    \item prend davantage de paramètres en compte ;
    \item fournit directement un facteur de stabilité ;
    \item est mieux adaptée aux projectiles modernes.
\end{itemize}

\section{Limites de la formule}

Malgré ses qualités, la formule reste une approximation.

Elle ne remplace pas :

\begin{itemize}
    \item les essais réels ;
    \item les mesures expérimentales ;
    \item les simulations avancées.
\end{itemize}

Les phénomènes dynamiques complexes ne sont pas entièrement représentés.

\section{Utilisation pratique}

La formule de Miller est largement utilisée :

\begin{itemize}
    \item par les tireurs ;
    \item par les rechargeurs ;
    \item par les concepteurs de projectiles.
\end{itemize}

Elle constitue souvent un excellent point de départ pour évaluer la
compatibilité entre un projectile et un canon.

\section{Lien avec la stabilité dynamique}

La formule de Miller évalue essentiellement la stabilité gyroscopique.

Cependant, un projectile peut être gyroscopiquement stable tout en présentant
des comportements dynamiques indésirables.

Pour comprendre ces phénomènes, il faut introduire la notion de stabilité
dynamique.

\section{À retenir}

\begin{aretenir}

\begin{itemize}
    \item La formule de Miller estime directement le facteur de stabilité
          gyroscopique.
    \item Elle est plus complète que la formule de Greenhill.
    \item La longueur du projectile joue un rôle majeur.
    \item Les conditions atmosphériques influencent le résultat.
    \item Une bonne stabilité gyroscopique ne garantit pas à elle seule un
          comportement dynamique parfait.
\end{itemize}

\end{aretenir}
